\section{研究综述}

\subsection{LLM 驱动的自动化数据分析研究进展}

LLM 驱动的数据分析系统近年来快速发展，推动了数据科学工作流从人工操作向自动化编排演进\cite{sun2025survey,zhao2024llm}。从研究趋势看，数据分析任务已不再被简单视为“让模型写一段代码”，而是逐步被重构为由任务理解、方案规划、代码执行、结果解释与反馈修正等多个环节组成的复合流程。代表性工作表明，多智能体机制能够将复杂数据分析过程拆分为“任务规划、代码执行、结果解释”等子环节，从而提高自动化完成率和整体稳定性。

DS-Agent\cite{guo2024dsagent} 通过案例推理（Case-Based Reasoning, CBR）从 Kaggle 历史任务库中检索相似案例，利用迁移策略指导 LLM 生成数据科学代码。该方法显著增强了任务求解的可复用性，但其核心依赖在于历史案例库的覆盖度，在跨领域或数据分布显著变化时可能失效。Data Interpreter\cite{hong2025data} 通过“动态规划—代码执行—反馈迭代”循环增强了代码执行的稳定性和正确率，是当前 LLM 数据科学 Agent 中执行层表现突出的代表性系统之一。二者共同说明，现有自动化分析系统已经能够在“给定分析目标后如何更稳地做出来”这一层面取得较大进展。

Data-to-Dashboard\cite{zhang2025data} 面向企业分析场景，通过多角色协同生成可视化仪表板，展示了 LLM Agent 在商业分析中的应用前景。此外，相关综述指出，统计与数据科学场景中的 LLM Agent 生态正在形成，系统能力正从“单模型问答”向“任务编排执行”演进\cite{sun2025survey}。这意味着，自动化数据分析研究正在从“回答分析问题”逐渐迈向“承担分析工作流”。

尽管如此，现有研究仍主要优化“执行层”，对“规划层”的研究不足。所谓执行层，关注的是给定任务后如何正确产出代码和结果；而规划层要解决的是“应该分析什么、为何这样分析、哪些变量需要控制”。在科研型数据分析中，规划层质量往往决定最终结论质量。若规划阶段缺乏数据证据，后续执行再稳定，也可能只是“稳定执行了不够好的方案”。因此，相较于执行正确率的持续提升，如何提升规划阶段的质量与针对性，正在成为自动化分析系统进一步突破的关键。

另一个共性问题是“历史案例依赖”与“当前数据脱钩”。DS-Agent 使用历史任务库强化迁移能力，但在跨领域或数据分布显著变化时，历史最优方案未必适配当前数据。Data Interpreter 的动态规划虽然在执行过程中根据代码反馈调整策略，但规划的起点仍是问题语义而非数据事实，系统并未在规划前对数据的变量分布、跨列关系和结构特征进行系统性探测。专利分析场景中这一问题更为突出，因为字段语义与技术背景强耦合，指标解释需要结合领域上下文，模板迁移风险较高。

从“规划是否感知数据”这一维度审视现有系统，可以将其分为两类：\textbf{数据无关规划}（data-blind planning）和\textbf{数据感知规划}（data-aware planning）。前者仅依据问题语义和历史经验生成分析方案，代表工作包括 DS-Agent 的案例检索和通用 LLM 的直接指令生成；后者则在规划前显式探测当前数据特征，将数据事实作为方案设计的输入约束。该划分有助于更清晰地识别现有自动化分析系统的能力边界，也为本文的研究切入点提供了概念基础。

现有系统中，Data-to-Dashboard\cite{zhang2025data} 通过多轮对话让用户补充数据信息，具有隐式的数据感知特征，但其感知过程依赖人工交互而非自动化提取，且仅面向可视化而非统计分析。综合来看，自动化、显式面向科研规划的数据感知机制仍相对不足，系统化方案较少，这也正是本文在自动化分析方向上的主要切入点。

\subsection{多智能体协作框架研究进展}

多智能体系统为 LLM 应用提供了结构化的协作范式。ReAct\cite{yao2023react} 提出将推理（Reasoning）与行动（Acting）在语言模型中交替执行，建立了 LLM Agent 的基本行为模式，即先思考、再执行，并根据环境反馈持续调整策略。Reflexion\cite{shinn2024reflexion} 进一步引入语言强化学习机制，使智能体能够从执行失败中学习并自我改进。这一范式为后续多智能体架构奠定了基础，也说明语言模型系统正在从“单轮响应器”向“可持续交互的任务执行体”演化。

在多智能体协作框架方面，AutoGen\cite{wu2023autogen} 提出基于多智能体对话的编程范式，支持灵活定义 Agent 角色、交互规则和终止条件，已在代码生成、数学推理等任务中展示了多 Agent 协作的优势。MetaGPT\cite{hong2024metagpt} 引入“标准作业程序”（SOP）概念，将软件工程的角色分工（产品经理、架构师、工程师）映射到 LLM Agent 协作流程中，通过结构化中间产物传递信息、减少幻觉。CAMEL\cite{li2023camel} 通过角色扮演和引导式对话实现两个 LLM Agent 之间的自主合作，探索了无人干预情况下的协作涌现能力。AgentBench\cite{liu2023agentbench} 提供了评估 LLM 作为智能体能力的标准化基准测试，涵盖多种环境和任务类型。这些工作共同表明，多智能体框架已经能够较好地支持角色分工、信息传递与任务拆解。

在工程实现层面，LangChain\cite{chase2022langchain} 提供了 LLM 应用的基础组件库，包括提示模板、工具调用和检索增强等能力；LangGraph\cite{langgraph2024} 在其基础上引入有向图状态机机制，支持循环、分支和条件路由，适合编排需要多步决策和反馈循环的复杂工作流。相较于 AutoGen 的对话驱动模式，LangGraph 的状态机模式在流程可控性和中间状态可追踪性方面具有优势\cite{wang2024survey,xi2023rise}。此外，工具学习（Tool Learning）研究表明，LLM 可以学习使用外部工具来增强自身能力\cite{schick2024toolformer,qin2023tool,mialon2023augmented}，这为构建具有复杂功能的智能体系统提供了理论基础。Voyager\cite{wang2023voyager} 和 Generative Agents\cite{park2023generative} 等工作则进一步展示了 LLM 智能体在开放环境中的自主学习和交互能力。由此可见，多智能体系统的工程基础设施已日趋成熟。

综合来看，当前多智能体框架已较好解决“如何让多个 Agent 协作”的基础工程问题，但在“Agent 应该做什么决策”这一更高层面上仍高度依赖提示词设计。对于数据分析场景而言，现有框架更多提供的是协作骨架，而非面向特定领域的决策机制，因此仍缺少数据感知和领域知识注入的内建能力。

\subsection{大语言模型在专利分析中的应用进展}

大语言模型与专利任务的结合已经从早期摘要生成逐步拓展到问答、语义检索、知识获取与辅助评审等更复杂的应用场景。与通用技术文本相比，专利文本具有句式冗长、术语密集、权利要求结构复杂、法律语言与技术语言交织等特点，这使其既适合发挥大语言模型在长文本理解与语义归纳上的优势，也对模型的领域适配能力提出了更高要求。EvoPat\cite{wang2024evopat} 通过多大语言模型智能体实现专利摘要、比对与评估流程，利用检索增强生成和子任务分解策略提升处理效果。KLIPA\cite{deng2025klipa} 引入知识图谱增强专利问答和信息获取，将结构化知识与非结构化文本检索相结合。与此同时，专利自然语言处理综述持续指出，专利文本处理仍是当前研究主流，涵盖分类、检索、实体识别、语义匹配等任务\cite{jiang2025natural}。总体来看，现有研究已经较充分证明，大语言模型在专利语义理解、信息抽取和文本级知识服务方面具有明显潜力。

这些进展证明了大语言模型在“专利文本理解”上的价值，但仍难直接替代“专利计量分析”。原因在于，文本任务的核心目标通常是从专利文献中识别、归纳和组织语义信息，而计量分析的核心目标则是围绕结构化专利数据识别变量关系、刻画时间演化、解释竞争格局并建立机制性判断。后者不仅要求系统理解文本含义，还要求其能够完成变量测量、指标构造、模型选择、混杂控制与稳健性检验等一整套统计分析链条\cite{basberg1987patents,trajtenberg1990penny,daim2006forecasting}。换言之，专利计量分析并不是“把专利文本读懂”就可以完成的任务，而是需要把文本、元数据、引文网络、国家与机构属性等多种信息统一转化为可执行的分析方案。该链条与文本问答的技术路径差异显著，也意味着面向专利分析的大语言模型系统不能仅停留在检索与回答层面，而必须进一步具备方案规划与分析执行能力。

已有研究也提示，专利分析正从单点文本任务向多维度融合任务演进，亟需将文本、结构字段、网络关系和统计建模统一到一个可执行框架中\cite{zhu2014patent,mejia2024patent}。例如，在技术路线识别、竞争态势评估、创新能力比较和知识扩散分析等典型问题中，研究者往往需要同时结合 IPC/CPC 分类、专利族信息、申请人属性、引证关系、时间序列特征以及文本语义信号，单一文本处理方法难以支撑此类综合性研究。专利可视化分析\cite{park2012visualization}和技术聚类方法\cite{lee2015technology}为理解技术演化模式提供了重要工具，但这些方法仍需要人工设计分析流程。大学专利与技术转移研究\cite{chen2004university}表明，专利数据分析在产学研合作中具有重要应用价值。由此可见，当前研究虽然已经在“专利文本可理解”“专利知识可检索”层面取得明显进展，但在“专利结构化数据能否被自动组织为高质量分析方案”这一更高层次问题上，仍缺少成熟的方法体系与系统实践。因此，面向专利结构化数据的“方案级自动化”仍然是关键空白。

\subsection{知识图谱增强大语言模型研究进展}

知识图谱增强大语言模型的核心目标是将外部结构知识引入生成过程，以提升一致性、可解释性和事实正确性\cite{pan2024unifying,liu2016knowledge}。与仅依赖参数记忆的语言模型相比，知识图谱能够以实体、关系和层级结构的形式提供更稳定的语义约束，使模型在处理复杂领域问题时不仅“知道更多”，而且“组织得更有结构”。GraphRAG\cite{edge2024graphrag} 通过“局部检索 + 全局社区结构”提高了复杂问题回答质量，展示了图结构在长上下文组织中的优势。检索增强生成（RAG）技术\cite{gao2023retrieval}为 LLM 提供了访问外部知识库的能力，有效缓解了知识过时和幻觉问题。相关综述进一步指出，知识约束在降低幻觉\cite{ji2023survey,zhang2023siren}和提升推理稳定性方面有明显潜力\cite{zhou2024hypothesis}。总体而言，知识图谱增强 LLM 已经从“给模型补充事实”逐步发展到“帮助模型组织推理路径”，这为复杂任务中的结构化决策奠定了方法基础。

但从应用目标看，当前多数知识图谱增强大语言模型工作仍聚焦检索问答，知识图谱主要作为“信息源”而非“决策器”。也就是说，图谱更多用于帮助模型找到正确知识、补全背景信息和减少事实性错误，而较少直接参与“应当选择什么分析路径”“哪些变量关系更值得优先验证”“哪类方法更适合当前任务”这类规划型决策。对于数据分析任务而言，方案质量不仅受知识完整性影响，更受数据可用性和方法可执行性影响。若知识约束过强而缺乏数据证据校准，系统可能陷入“结构正确但数据适配不足”的状态，即理论上看起来合理、但落到当前数据集上却难以执行或解释力有限。已有研究表明，LLM 作为知识库在语义理解任务中具有良好表现\cite{kim2023language}，但在数据分析场景中如何有效利用 LLM 的知识能力仍需探索。此外，领域适应研究\cite{wu2024brief}提示，模型在不同数据分布下的泛化能力是影响分析方案迁移性的持续挑战。这说明，单纯把知识图谱“接到”语言模型上，并不自然等于形成了高质量分析系统，关键仍在于知识如何与数据证据、方法选择和任务目标发生协同。

在更广泛的大语言模型科学发现方向，AI Scientist v2\cite{yamada2025aiscientist} 实现了从假设生成到实验执行到论文撰写的全自动科研流程。EMNLP 2025 的大语言模型科学发现综述\cite{zheng2025automation} 定义了三个自主水平：Level 1（LLM as Tool）、Level 2（LLM as Analyst）、Level 3（LLM as Scientist）。这些工作表明大语言模型正在从辅助工具向自主分析者演进，也意味着未来智能系统不再只是回答问题，而是有可能参与问题定义、方案设计、实验验证乃至知识生产的完整流程。然而，从当前 Level 2 系统的发展状态来看，它们往往已经具备一定的任务分解、代码执行和结果总结能力，但通常未将“规划前数据感知”和“领域知识图谱约束”同时作为核心设计组件。前者决定系统能否真正理解“当前数据能够支持什么样的分析”，后者决定系统能否在领域约束下形成更稳定、更可解释的分析路径。正是在这一意义上，知识图谱增强 LLM 的研究虽然为本文提供了重要启发，但仍未直接解决面向专利结构化数据的规划型分析问题。

为避免仅凭文字叙述给出研究定位，表\ref{tab:related-comparison} 按本文最关心的五个能力维度对代表性工作进行归纳。表中判断依据均来自文献中显式报告的系统能力，而非潜在可扩展能力；其中“弱/部分”表示该能力存在，但并非系统的主要设计目标或未形成稳定工作流。

\begin{table}[htbp]
\centering
\caption{相关代表性工作与本文的能力维度对比}
\label{tab:related-comparison}
\resizebox{\textwidth}{!}{
\begin{tabular}{@{}cccccc@{}}
\toprule
工作 & 主要场景 & 规划前数据感知 & 领域知识约束 & 面向专利结构化分析 & 反馈迭代 \\ \midrule
DS-Agent & 通用数据科学 & 否（依赖历史案例） & 否 & 否 & 弱 \\
Data Interpreter & 通用数据科学 & 弱（执行中反馈） & 否 & 否 & 是 \\
Data-to-Dashboard & 企业分析可视化 & 弱（依赖用户补充） & 否 & 否 & 是 \\
EvoPat / KLIPA & 专利文本任务 & 否或弱 & 是（检索/图谱） & 否（偏文本） & 部分 \\ \midrule
\textbf{本文} & \textbf{专利结构化计量分析} & \textbf{是（自动预览）} & \textbf{是（双知识图谱）} & \textbf{是} & \textbf{是} \\ \bottomrule
\end{tabular}
}
\end{table}

\subsection{研究现状总结与本文定位}

综合前述综述及表\ref{tab:related-comparison}可以看出，当前相关研究虽然分别在自动化数据分析、多智能体协作、专利文本处理以及知识图谱增强大语言模型等方向上取得了显著进展，但这些进展大多仍停留在各自子问题的局部优化层面，尚未形成一个同时覆盖“规划前数据感知、领域知识约束、专利结构化分析、反馈迭代”四个维度的统一框架。换言之，现有工作已经为本文提供了重要的方法启发和技术积累，但在规划层能力、领域适配能力以及知识与数据协同机制方面仍存在明显缺口。进一步归纳来看，这些不足主要体现在以下三个方面，且分别对应本文的一项核心研究内容：


（1） 在大语言模型自动化分析方向，执行优化多、规划优化少，缺乏“规划前数据感知”机制。现有系统普遍更关注任务执行是否成功、代码是否可运行、结果是否可输出，而较少关注分析方案在生成之前是否已经充分理解当前数据的分布特征、结构关系和潜在约束。对此，本文通过数据预览、图结构预览与洞察生成三层输入，将数据证据前移至规划阶段，使方案设计建立在对当前数据集的显式感知基础之上，从而提升规划结果的针对性与可执行性。

（2） 在大语言模型专利应用方向，文本处理多、结构化计量少，缺少端到端方案生成系统。现有研究已经证明了大语言模型在专利摘要、检索、问答和知识抽取中的有效性，但对于专利计量分析这类依赖变量构造、模型选择和机制解释的结构化任务，尚未形成成熟的自动化系统路径。对此，本文构建面向专利结构化数据的四智能体系统，覆盖从问题输入、方案规划、方法规格化、代码执行到结果输出的完整流程，以填补专利分析从“文本理解自动化”走向“分析方案自动化”的关键空白。

（3） 在知识图谱增强大语言模型方向，检索问答多、分析引导少，缺乏“知识约束 + 数据证据”协同框架。现有知识图谱增强方法往往把图谱作为外部知识源来提升回答质量，但较少进一步探讨图谱如何参与分析路径规划、变量关系筛选和方法匹配决策。对此，本文提出双知识图谱协同机制，将因果图谱与方法图谱分别作用于研究假设空间与分析执行空间，并通过实验验证知识约束发挥正向作用的边界条件，从而更系统地回答“知识如何真正服务于分析规划”这一问题。


基于上述分析，本文的研究定位可以概括为：面向专利结构化数据分析场景，构建一个融合数据感知、双知识图谱约束与多智能体迭代优化机制的自动化分析系统。上述定位也使第二章与后续章节形成清晰映射：第二章提出问题缺口与理论动机，第三章和第四章分别给出系统设计方案与原型实现，第五章和第六章通过实验评估与应用案例验证系统有效性，第七章总结研究贡献并讨论其适用边界与未来拓展方向。由此，第二章不仅完成了相关工作的回顾，也为全文后续章节的展开建立了明确的问题导向和方法起点。
